\newsec{Conclusioni}

Con questa esperienza abbiamo caratterizzato lo spettro di emissione di un tubo radiogeno, studiando anche l'interazione dei raggi X emessi con la materia.

All'inizio abbiamo calibrato il contatore Geiger-Muller fissando la tensione di alimentazione $ V_\text{alimentazione} = 506 \unit{V} $, e fornendo anche una stima del suo tempo morto:
\begin{equation}
\tau_\text{dead} = (187.8 \pm 1.5) \unit{\micro s}
\end{equation}
compatibilmente con le specifiche dello strumento.
Nello specifico abbiamo anche fissato le condizioni di lavoro ottimali del tubo radiogeno, scegliendo correnti di alimentazione del filamento $ i \leq 0.3 \unit{\micro A} $ per rimanere nella regione di risposta lineare del contatore.

Poi, abbiamo studiato la validità della legge di Lambert-Beer: il modello non si è adattato perfettamente ai dati a causa della natura non monocromatica del fascio di raggi X emessi.
Comunque, siamo riusciti a distinguere una componente a bassa energia associata alla radiazione caratteristica del rame con coefficiente di assorbimento:
\begin{equation}
\frac{\mu_b}{\rho} = (254 \pm 13) \unit{cm^2 / g}
\end{equation}
e una a energia maggiore associata al picco di bremsstrahlung:
\begin{equation}
\frac{\mu_a}{\rho} = (28 \pm 5) \unit{cm^2 / g}.
\end{equation}

Dopo, abbiamo avuto modo di studiare la diffrazione dello spettro dei raggi X su diversi cristalli, usando la legge di Bragg (\ref{eq:bragg}).

Nel caso del \ce{NaCl} abbiamo usato come dato la distanza reticolare del cristallo per stimare le lunghezze d'onda dei raggi X caratteristici del rame $K_\alpha$ e $K_\beta$:
\begin{equation}
\lambda_1 = (1.33 \pm 0.04) \unit{\angstrom} \qquad \lambda_2 = (1.48 \pm0.03) \unit{\angstrom}
\end{equation}
e già da queste stime abbiamo notato la presenza di un errore sistematico di $ 0.7 \unit{\degree} $ (scala $ \theta $, quindi per la scala di lettura è di $ 1.4 \unit{\degree} $) nella lettura degli angoli di diffrazione.

Per controllare già in fase di presa dati la presenza di un sistematico sull'azzeramento della scala con cui abbiamo misurato gli angoli avremmo potuto prendere delle misure simmetriche e opposte.
Ottenendo conteggi diversi ad angoli ``uguali'' avremmo potuto azzerare correttamente la scala, senza dover correggere le stime a posteriori.

Poi abbiamo studiato l'effetto di un filtro di \ce{Ni} sullo spettro.
Come nel caso dell'analisi fatta sull'assorbimento da parte dell'alluminio al variare dello spessore, abbiamo osservato delle soppressioni diverse per le due righe di emissione (\cref{fig:nichelmu} e \cref{fig:Ni}).\footnote{Preparando il fascio di raggi X con un filtro di nichel avremmo potuto verificare più facilmente la legge di Lambert-Beer: la condizione di monocromaticità sarebbe stata migliore, anche se non ancora perfetta.}
Solo il picco relativo a $K_\alpha$ era ben riconoscibile, e ne abbiamo stimato la lunghezza d'onda come:
\begin{equation}
\lambda_{2\ce{Ni}} = (1.48 \pm 0.03) \unit{\angstrom}
\end{equation}
a meno del sistematico sulla lettura degli angoli.

Infine, per i cristalli di $\ce{LiF}$ e $\ce{KCl}$ abbiamo misurato la loro distanza interatomica:
\begin{equation}
2d_{\ce{LiF}} = (4.14 \pm 0.07) \unit{\angstrom}\qquad 2d_{\ce{KCl}} = (6.5 \pm 0.2) \unit{\angstrom}
\end{equation}
e correggendo per il sistematico:
\begin{equation}
  2d_{\ce{LiF}}' = (4.02 \pm 0.09) \unit{\angstrom} \qquad 2d_{\ce{KCl}}' = (6.2 \pm 0.2) \unit{\angstrom}
\end{equation}
ottenendo delle stime compatibili con i valori nominali.
